% 集合标签类 H_C 的两种学习规则（教学示例，n=4，输入均匀分布在 {1,...,8}）。
% 上：目标 h_A，A={2,5}。样本输入 {1,5,6,7}，x=5 的标签就是集合 {2,5}，
%     读出标签的学习器恢复 h_A，风险 0。
% 下：目标 h_∅（全为 ★）。样本输入 {1,4,6,7}，标签全为 ★。
%     坏 ERM 输出 h_{A_S}，A_S={1..8}\{1,4,6,7}={2,3,5,8}；样本上全对，
%     但在 2,3,5,8 上错，风险 4/8=1/2。
\begingroup
\begin{tikzpicture}[figure picture,x=1mm,y=1mm]
 \tikzset{ltsixpt/.style={circle,draw=FigureStroke,line width=.6pt,fill=white,minimum size=5.4mm,inner sep=0pt,font=\figurefont\scriptsize},
   ltsixseen/.style={ltsixpt,fill=FigureBlueFill,draw=FigureBlue,line width=.9pt},
   ltsixlab/.style={font=\figurefont\scriptsize,text=FigureInk},
   ltsixbad/.style={ltsixlab,text=FigureRed}}
 \def\dx{9.5}
 % ---------- top panel
 \node[figure heading,anchor=west] at (0,47) {目标 $h_A$，$A=\{2,5\}$：样本含点 5};
 \foreach \i in {1,...,8}{
   \ifnum\i=1 \def\st{ltsixseen}\else\ifnum\i=5 \def\st{ltsixseen}\else\ifnum\i=6 \def\st{ltsixseen}\else\ifnum\i=7 \def\st{ltsixseen}\else\def\st{ltsixpt}\fi\fi\fi\fi
   \node[\st] at ({22+\dx*(\i-1)},38) {\i};}
 \node[ltsixlab,anchor=east] at (17,32) {真实标签};
 \foreach \i in {1,...,8}{
   \ifnum\i=2 \def\l{$\{2,5\}$}\else\ifnum\i=5 \def\l{$\{2,5\}$}\else\def\l{$\star$}\fi\fi
   \node[ltsixlab] at ({22+\dx*(\i-1)},32) {\l};}
 \node[ltsixlab,anchor=east] at (17,26) {读出标签};
 \foreach \i in {1,...,8}{
   \ifnum\i=2 \def\l{$\{2,5\}$}\else\ifnum\i=5 \def\l{$\{2,5\}$}\else\def\l{$\star$}\fi\fi
   \node[ltsixlab] at ({22+\dx*(\i-1)},26) {\l};}
 \node[figure note,anchor=west] at (95,26) {风险 $0$};
 % ---------- bottom panel
 \node[figure heading,anchor=west] at (0,15) {目标 $h_\varnothing$：样本标签全为 $\star$};
 \foreach \i in {1,...,8}{
   \ifnum\i=1 \def\st{ltsixseen}\else\ifnum\i=4 \def\st{ltsixseen}\else\ifnum\i=6 \def\st{ltsixseen}\else\ifnum\i=7 \def\st{ltsixseen}\else\def\st{ltsixpt}\fi\fi\fi\fi
   \node[\st] at ({22+\dx*(\i-1)},6) {\i};}
 \node[ltsixlab,anchor=east] at (17,0) {真实标签};
 \foreach \i in {1,...,8}{\node[ltsixlab] at ({22+\dx*(\i-1)},0) {$\star$};}
 \node[ltsixlab,anchor=east] at (17,-6) {读出标签};
 \foreach \i in {1,...,8}{\node[ltsixlab] at ({22+\dx*(\i-1)},-6) {$\star$};}
 \node[figure note,anchor=west] at (95,-6) {风险 $0$};
 \node[ltsixlab,anchor=east] at (17,-12) {坏ERM};
 \foreach \i in {1,...,8}{
   \ifnum\i=1 \def\l{$\star$}\def\s{ltsixlab}\else\ifnum\i=4 \def\l{$\star$}\def\s{ltsixlab}\else\ifnum\i=6 \def\l{$\star$}\def\s{ltsixlab}\else\ifnum\i=7 \def\l{$\star$}\def\s{ltsixlab}\else\def\l{$A_S$}\def\s{ltsixbad}\fi\fi\fi\fi
   \node[\s] at ({22+\dx*(\i-1)},-12) {\l};}
 \node[figure note,anchor=west] at (95,-12) {风险 $1/2$};
 \node[figure note,anchor=west] at (0,-19) {蓝色圆为样本中的输入；$A_S=\{2,3,5,8\}$ 恰是未抽到的点。};
\end{tikzpicture}
\endgroup
